\chapter{Contraceptive intra-uterine device (IUD)}
\index{Contraceptive intra-uterine device (IUD)}

\section*{Definition and Overview}
A contraceptive intra-uterine device (IUD) is a small, T-shaped contraceptive device inserted into the uterus by a trained medical professional to provide long-acting reversible contraception (LARC). There are two primary types of IUDs available: the copper IUD (Cu-IUD), which is hormone-free and effective for up to 5 to 10 years, and the hormonal levonorgestrel-releasing IUD (LNG-IUD), which is effective for 3 to 8 years depending on the specific model (e.g., Mirena, Kyleena). Both options are exceptionally reliable and represent some of the most effective birth control methods available. In addition to contraception, LNG-IUDs are widely utilised in clinical practice as a primary therapeutic management option for heavy menstrual bleeding, dysmenorrhoea, and endometrial protection during hormone replacement therapy.

\section*{Epidemiology}
Use of IUDs has increased in many regions. Copper and hormonal IUDs are recommended contraceptive options for people of reproductive age, including adolescents and people who have not previously been pregnant, when clinically suitable. Availability, duration, cost, and funding vary by country and by device. Because the method does not depend on daily use, its typical-use effectiveness is close to its clinical-trial effectiveness.

\section*{Aetiology and Pathophysiology}
The physiological mechanisms through which IUDs prevent pregnancy differ between the copper and hormonal variants, although both primarily function by preventing fertilisation.
\begin{itemize}
    \item \textbf{Copper-induced sterile inflammation (Cu-IUD):} The release of copper ions into the uterine cavity stimulates a local inflammatory reaction that is highly toxic to sperm and ova, impairing sperm motility, viability, and fertilization capacity.
    \item \textbf{Cervical mucus thickening (LNG-IUD):} Levonorgestrel causes the cervical mucus to become thick, sticky, and impermeable to sperm, blocking their entry into the upper reproductive tract.
    \item \textbf{Endometrial suppression (LNG-IUD):} Locally released progestogen induces profound endometrial atrophy, rendering the uterine lining thin, inactive, and hostile to implantation.
    \item \textbf{Foreign body response (Both):} The physical presence of a foreign object in the uterine cavity induces a mild, localized cellular response that interferes with fertilization and prevents blastocyst implantation.
    \item \textbf{Anovulation (LNG-IUD):} While most women continue to ovulate normally on an LNG-IUD, partial ovarian follicular suppression occurs in some individuals, particularly in the first year.
\end{itemize}

\section*{Clinical Features}
The clinical presentation of patients with an IUD is characterised by the method-specific alterations to their menstrual cycles and physical presence of the device.
\begin{itemize}
    \item \textbf{Bleeding changes (Copper IUD):} Menstrual blood loss often increases by 20\% to 50\%, accompanied by longer menstrual periods and an increase in cramping or dysmenorrhoea, especially during the first 3 to 6 months.
    \item \textbf{Bleeding changes (Hormonal IUD):} Irregular spotting or light, unpredictable bleeding is common for the first 3 to 6 months. Thereafter, overall blood loss is dramatically reduced, and many users experience oligomenorrhoea or complete amenorrhoea (up to 60\% of Mirena users by 12 months).
    \item \textbf{Post-insertion cramping:} Mild to moderate lower abdominal cramping and backache are common immediately following insertion, resolving within a few days.
    \item \textbf{String presence:} The thin monofilament strings extend through the cervical canal and are palpable at the external os, serving as a marker of proper placement.
    \item \textbf{Hormonal side effects (LNG-IUD):} Some women may experience transient symptoms such as breast tenderness, mild acne, mood changes, or functional ovarian cysts.
\end{itemize}

\section*{Differential Diagnosis}
When evaluating a patient with an IUD presenting with pelvic pain, abnormal bleeding, or pregnancy symptoms, several clinical conditions must be differentiated.
\begin{itemize}
    \item \textbf{Ectopic pregnancy:} Although the absolute risk of pregnancy is extremely low, if conception does occur with an IUD in situ, a high proportion (up to 50\%) of these pregnancies are ectopic.
    \item \textbf{IUD expulsion:} Partial or complete expulsion of the device from the uterine cavity must be considered if strings are missing, or if the patient experiences new-onset cramping and unexpected bleeding.
    \item \textbf{Uterine perforation:} A rare insertion complication (occurring in 1 to 2 per 1,000 insertions) where the device passes through the myometrium, presenting with severe acute pain or asymptomatic loss of strings.
    \item \textbf{Pelvic inflammatory disease (PID):} Bacterial infection of the upper genital tract, which may occur more frequently in the first 20 days post-insertion if pre-existing STIs were present.
    \item \textbf{Gynaecological pathology:} Bleeding or pain may be due to unrelated conditions such as cervical ectropion, uterine fibroids, endometrial polyps, or sexually transmitted infections.
\end{itemize}

\section*{When to Seek Medical Care}
Patients should consult a general practitioner or sexual health clinic if they cannot feel their IUD strings, feel the hard plastic of the device at the cervix, experience abnormal discharge, or have severe, persistent cramping.

\textbf{Contact your local emergency services for:}
\begin{itemize}
    \item Severe, worsening lower abdominal or pelvic pain that does not respond to simple analgesics
    \item High fever, chills, and foul-smelling vaginal discharge (signs of pelvic infection)
    \item Heavy, uncontrolled vaginal bleeding soaking through multiple pads or tampons per hour
    \item Sudden dizziness, fainting, or severe shoulder-tip pain (possible signs of a ruptured ectopic pregnancy)
    \item Persistent vomiting, severe bloating, or signs of systemic sepsis
\end{itemize}

\section*{Diagnosis and Investigations}
Confirming the correct positioning of an IUD and ruling out complications requires specific diagnostic evaluations.
\begin{itemize}
    \item \textbf{Pelvic examination:} Speculum examination to visualize the IUD strings at the external cervical os, and bimanual palpation to check for uterine tenderness.
    \item \textbf{Pelvic ultrasound:} Transvaginal ultrasound is the gold standard imaging modality to confirm the precise intrauterine location of the IUD and rule out displacement, partial expulsion, or myometrial penetration.
    \item \textbf{Pregnancy test:} Urine or serum beta-hCG to exclude pregnancy in patients presenting with amenorrhoea (especially with a copper IUD) or pelvic pain.
    \item \textbf{STI screening:} Swabs for Chlamydia trachomatis and Neisseria gonorrhoeae, ideally performed prior to or at the time of insertion.
    \item \textbf{Abdominal X-ray:} Performed if the IUD is not seen on ultrasound and perforation into the peritoneal cavity is suspected, as copper and hormonal IUDs are radiopaque.
\end{itemize}

\section*{Management}
The clinical management of IUDs involves comprehensive patient selection, structured insertion procedures, and follow-up care.
\begin{itemize}
    \item \textbf{Pre-insertion assessment:} Detailed medical history to check for contraindications (e.g., uterine anomalies, active pelvic infection, unexplained bleeding, or breast cancer for hormonal IUDs).
    \item \textbf{Insertion procedure:} Performed under sterile conditions. It involves measuring uterine depth (sounding) and passing the device through the cervical canal into the fundus. Local anaesthesia, such as a cervical block or topical lignocaine gel, is often utilised to minimise discomfort.
    \item \textbf{Pain and spasm management:} Suggesting pre-procedure NSAIDs (e.g., ibuprofen) to reduce uterine cramping, and ensuring the patient has a resting period immediately post-insertion.
    \item \textbf{Follow-up appointment:} Typically scheduled 4 to 6 weeks post-insertion to check the strings, assess tolerance, and rule out early expulsion or infection.
    \item \textbf{Removal and replacement:} The device can be easily removed by pulling the strings with forceps. Immediate replacement can be performed during the same clinic visit if ongoing contraception is desired.
\end{itemize}

\section*{Complications}
While IUDs are highly safe, some patients may experience procedural or long-term complications.
\begin{itemize}
    \item \textbf{Expulsion:} The uterus may spontaneously expel the device, most commonly during the first few months or during menstruation (occurs in 5\% of users).
    \item \textbf{Uterine perforation:} A serious insertion risk where the device penetrates the uterine wall, requiring laparoscopic surgical retrieval.
    \item \textbf{Pelvic infection:} A slight increase in the risk of PID during the first few weeks after insertion, usually associated with pre-existing untreated cervical infection.
    \item \textbf{Menstrual disturbance:} Heavy bleeding and dysmenorrhoea with copper IUDs, or persistent irregular spotting and unpredictable bleeding with hormonal IUDs.
    \item \textbf{Lost strings:} The strings may retract into the cervical canal or uterine cavity, requiring ultrasound location and specialized instruments for retrieval or removal.
\end{itemize}

\section*{Prognosis and Outcomes}
The long-term prognosis for IUD users is exceptional. The contraceptive failure rate is less than 0.2\% for hormonal IUDs and 0.8\% for copper IUDs over the first year of use, which is comparable to surgical sterilization. Once the device is removed, the return to the patient’s baseline fertility is rapid, with no negative long-term impacts on future reproductive health or pregnancy outcomes. Hormonal IUDs also provide significant non-contraceptive benefits, reducing menstrual blood loss by up to 90\% and often resolving chronic pelvic pain associated with adenomyosis or endometriosis.

\section*{Prevention and Self-Care}
Patients can engage in self-care strategies to monitor their device and ensure optimal comfort.
\begin{itemize}
    \item \textbf{Perform monthly string checks:} Gently feel for the IUD strings at the top of the vagina after each period to confirm the device remains in place.
    \item \textbf{Avoid pulling strings:} Never pull on the strings, as this can displace or expel the IUD.
    \item \textbf{Prepare for insertion:} Take recommended oral pain relief (such as ibuprofen) 1 hour before the insertion appointment as advised by the clinician.
    \item \textbf{Observe pelvic rest:} Avoid inserting anything into the vagina (including tampons, cups, and sexual intercourse) for 48 hours after insertion to minimise infection risk.
    \item \textbf{Track symptoms:} Maintain a menstrual diary to monitor changes in bleeding patterns and pelvic pain, particularly during the initial adaptation period.
\end{itemize}

\section*{Sources and Further Reading}
The following reputable organisations provide further information on IUDs:
\begin{itemize}
    \item Family Planning Alliance Australia. \textit{Intrauterine devices (IUDs) clinical guide.} https://www.fpalliance.org.au/
    \item Healthdirect Australia. \textit{Intrauterine device (IUD) overview.} https://www.healthdirect.gov.au/iud-intrauterine-device
    \item Royal Australian and New Zealand College of Obstetricians and Gynaecologists (RANZCOG). \textit{Patient resource on IUDs.} https://ranzcog.edu.au/
    \item NHS. \textit{Intrauterine device (IUD) and intrauterine system (IUS) guidelines.} https://www.nhs.uk/contraception/methods/iud/
    \item Mayo Clinic. \textit{Mirena (hormonal IUD) and copper IUD information.} https://www.mayoclinic.org/tests-procedures/mirena/about/pac-20391354
\end{itemize}
